\documentclass[11pt]{article}
\usepackage[T1]{fontenc}
\usepackage[utf8]{inputenc}
\usepackage{lmodern}
\usepackage{amsmath,amssymb,booktabs}
\usepackage[a4paper,margin=1in]{geometry}

\title{Unified Quartic--Sextic Axis-Representation Transform Program}
\author{}
\date{}

\begin{document}
\maketitle

\section*{Scope}
A single Python GUI program was implemented to replace the previous quartic-only Yamada app. The new program includes:
\begin{itemize}
\item quartic Yamada transform in A reduction,
\item quartic Yamada-like transform in S reduction,
\item quartic tensor algorithm (Watson $\rightarrow \tau \rightarrow$ permutation $\rightarrow$ Watson),
\item sextic tensor algorithm (Watson $\rightarrow \phi \rightarrow$ permutation $\rightarrow$ Watson).
\end{itemize}

Representations are $I,II,III$. Reduction is preserved (input reduction equals output reduction).

\section*{Implemented Quartic Models}
For quartic constants
\[
\mathbf D_A=(A_J,A_{JK},A_K,d_J,d_K)^\mathsf T,
\qquad
\mathbf D_S=(D_J,D_{JK},D_K,d_1,d_2)^\mathsf T,
\]
the reference transform is
\[
\mathbf D_{R'}=Q_{R'}^{-1}Q_R\,\mathbf D_R,
\]
where $Q_R=Q_A(R)$ for A reduction and $Q_R=Q_S(R)$ for S reduction.

For S reduction, the explicit implemented matrices are derived from S-Hamiltonian coefficient matching and use the corrected definition
\[
T_1=\frac12\left[\mathrm{coeff}(J_a^2J_b^2)+\mathrm{coeff}(J_a^2J_c^2)+\mathrm{coeff}(J_b^2J_c^2)\right].
\]

The quartic tensor algorithm is implemented as
\[
\mathbf D \xrightarrow{\,M_4^+\,} \tau \xrightarrow{\,P\,} \tau' \xrightarrow{\,M_4'\,} \mathbf D',
\]
with pseudoinverse $M_4^+$ and representation permutation matrix $P$.

\section*{Implemented Sextic Tensor Model}
Sextic Watson vector (A notation):
\[
\mathbf H=(H_J,H_{JK},H_{KJ},H_K,h_1,h_2,h_3)^\mathsf T.
\]
Sextic tensor vector:
\[
\phi=(\phi_{aaaaaa},\phi_{bbbbbb},\phi_{cccccc},\phi_{aaaabb},\phi_{aaaacc},\phi_{bbbbaa},\phi_{bbbbcc},\phi_{ccccaa},\phi_{ccccbb},\phi_{aabbcc})^\mathsf T.
\]
Mapping and transform:
\[
\mathbf H=M_6\phi,
\qquad
\phi=M_6^+\mathbf H,
\qquad
\mathbf H' = M_6' P_6 \phi.
\]
Here $P_6$ is the sextic index-permutation operator for representation changes ($I\to II$, $I\to III$, $II\to III$). Rotational constants are cyclically reordered along the same axis change.

\section*{Validation Results}
\subsection*{Quartic: reference vs tensor}
Using COF$_2$ constants (Carpenter), the implemented test script gives machine-precision agreement:
\begin{itemize}
\item $I_r\to II_r$: $\max|\Delta| = 1.776\times 10^{-15}$
\item $I_r\to III_r$: $\max|\Delta| = 2.776\times 10^{-17}$
\item $II_r\to III_r$: $\max|\Delta| = 2.776\times 10^{-17}$
\end{itemize}
and matrix identity check
\[
M(A',B',C')\,P\,M^+(A,B,C)=Q_j^{-1}Q_i
\]
returns zero (within floating-point precision) for all three pairs.

\subsection*{Sextic: round-trip stability}
For random sextic constants, round-trip tests
\[
I_r\to II_r\to I_r,
\quad
I_r\to III_r\to I_r,
\quad
II_r\to III_r\to II_r
\]
give
\[
\max\left|\mathbf H_{\mathrm{back}}-\mathbf H_{\mathrm{orig}}\right|=0.0
\]
in the executed numerical runs (machine precision).

\section*{Delivered Files}
\begin{itemize}
\item \texttt{yamada\_gui.py}: unified GUI program (quartic + sextic)
\item \texttt{test\_tensor\_algorithm.py}: quartic tensor/reference validation
\item \texttt{test\_tensor\_sextic\_algorithm.py}: sextic tensor round-trip validation
\item \texttt{build\_app.sh}: app builder updated to produce \texttt{CentrifugalTransform.app}
\end{itemize}

\end{document}
